% =====================================================================
\section{결론}
\label{sec:concl}

본 연구는 이미 통지된 외부 트럭 작업의 야드 블록과 도착시각을 운영 중에 함께
재검토하는 강화학습 아키텍처를 제안했다. 이 구조는 분산 예약 협상~\cite{ref6}과 야드
크레인 운영~\cite{ref7,ref8}을 하나의 시공간 배치 문제로 연결한다. 재배치의 영향은
뒤따르는 트럭, 크레인과 선박 작업에서 나중에 나타난다. 따라서 각 정책은 즉시
이동비보다 고정 지평에서 측정한 반사실 비용 차이를 학습한다. 이는 다중 행위자 반사실
학습~\cite{ref10,ref11}을 운영 의사결정에 적용하는 방법이다. 계산량이 큰 반사실 분기는
오프라인 학습에만 사용하므로, 운영 경로에는 작은 비용모형과 제약 검사만 남는다.

\ref{sec:env}절의 합성환경에서 제안 구조는 관측부터 제안, 수락, 확정과 실행까지 시공간
재배치의 전체 경로를 수행했다. 비용 감소는 주로 혼잡 조건의 도착시각 조정에서
나타났다. 즉 실시간 혼잡 정보가 모니터링에 그치지 않고 배치 행동으로 이어졌다. 같은 행동 범위를 가진 규칙과 비교해도 총비용은
학습 정책이 더 낮았다. 다만 시간 행동만 비교하면 규칙이 더 많은 날에서 낮았고, 학습
정책은 혼잡한 날의 큰 감소로 총액에서 앞섰다. 따라서 결과는 승리한 날짜 수와 총
감소액을 함께 보아야 한다.

\subsection{연구 범위와 확장 방향}
\label{sec:scope}
본 연구는 문헌에서 보고된 터미널 규모와 운영 변동 범위를 바탕으로 구성한 합성환경에서
제안 아키텍처의 작동 가능성과 시공간 재배치의 비용 효과를 검증하였다. 따라서 현재
결과는 특정 터미널에 대한 직접적인 성능 추정보다는, 운영 중 관측되는 혼잡 정보를
이용해 기존 배정을 재검토하는 구조가 어떠한 조건에서 유효한지를 보여 주는 실험적
근거로 해석할 수 있다. 향후 실제 터미널의 도착분포, 수요 구성과 비용 자료를 적용하면
서로 다른 운영환경에 대한 일반화 범위를 보다 구체적으로 평가할 수 있다.
특히 본 실험에서는 비용 감소가 높은 수요의 혼잡 조건에 집중되었다. 이는 재배치 정책을
상시 적용하기보다 혼잡 수준에 따라 선택적으로 작동시키는 방향으로 확장할 가능성을
보여 준다. 향후 다양한 수요 패턴과 혼잡 빈도를 반복적으로 평가하면 이러한 선택적 개입
전략의 적용 조건과 학습 효과를 보다 안정적으로 규명할 수 있다.
운영 측면에서는 실제 예약시스템의 시간대별 정원, 운송사와 터미널 간 예약 변경 절차,
그리고 다양한 크레인 작업순서 규칙을 함께 고려하는 확장이 필요하다. 본 연구에서 사용한
개방형 시간대 정원은 도착시각 조정이 제공할 수 있는 잠재적 효과를 평가하기 위한
설정이며, 실제 예약 정원을 적용하면 정책의 실행 가능한 행동 범위를 보다 현실적으로
정의할 수 있다. 또한 제안과 수락을 분리한 현재 구조는 향후 수락 거부권의 유무와 다양한
협상 구조를 비교함으로써 각 정책 구성요소의 역할을 보다 세밀하게 분석할 수 있다.
방법론적으로는 3시간의 반사실 평가 지평과 비용 계수가 정책 행동의 상대적 가치에 미치는
영향을 추가로 분석할 수 있다. 특히 현재 비용구조에서는 트럭 대기비의 비중이 높아
도착시각 조정의 효과가 크게 나타났으므로, 다양한 비용구조와 평가 지평에 대한 민감도
분석은 공간 및 시간 재배치가 유효한 운영 조건을 규명하는 데 도움이 될 것이다.
이러한 확장은 본 연구에서 제시한 관측--제안--수락--확정--실행--반사실 학습의 기본
구조를 변경하기보다, 이를 실제 운영조건과 다양한 수요환경으로 확장하는 과정에 해당한다.
따라서 본 연구는 실시간 혼잡 정보를 단순 모니터링에 머무르지 않고 설명 가능한 자원 배치
결정으로 연결하기 위한 하나의 기본 아키텍처를 제시하며, 향후 실제 터미널 자료와 운영
제약을 결합한 검증으로 발전할 수 있다.
